% =====================================================================
\lezione{16/11/2026}{NIS2 (II): governance, misure, notifiche, vigilanza, sanzioni; caso Microsoft 365}
  {Gianluca Gilardi (lezione online); breve intervento di Giuseppe Vaciago}
  {strategia nazionale e organismi UE (Cooperation Group, EU-CyCLONe); responsabilità personale
   dell'organo di gestione; misure proporzionate e perimetro NIS2; divulgazione coordinata delle
   vulnerabilità; incidente e incidente significativo; orologi 24 h / 72 h / 1 mese e orologio GDPR;
   vigilanza \emph{ex ante} / \emph{ex post}; esecuzione e sanzioni (triplicate per mancata
   registrazione); \emph{ne bis in idem} con il GDPR; caso del disservizio Microsoft 365; presentazione
   del laboratorio~5 (EnergiaTech)}
\label{lez:2026-11-16}
% =====================================================================

\begin{nota}[Fonte]
Trascrizione automatica della registrazione (lezione online; la lezione inizia dopo circa 8 minuti di
silenzio e dura circa 80 minuti), revisionata. Slide: il docente proietta le slide NIS2 aggiornate
(``NIS2\_PoliTo\_2026-27\_v3'', slide 17--40, con il caso Microsoft~365) e la presentazione del
laboratorio~5. Le slide sono in \emph{material} (\emph{08.~NIS2\_PoliTo\_2026-27\_v3}); il materiale del
laboratorio è nel Capitolo~\ref{lab:05}.
\end{nota}

\section{Sintesi}
Seconda e ultima lezione sulla NIS2. Il docente completa il quadro istituzionale (strategia nazionale,
Cooperation Group, EU-CyCLONe), poi passa agli obblighi: approvazione delle misure e
\textbf{responsabilità personale} dell'organo di gestione, misure \textbf{proporzionate} da valutare
solo sul \textbf{perimetro NIS2} del soggetto, notifica degli \textbf{incidenti significativi} con i suoi
orologi (distinti da quello del GDPR), vigilanza asimmetrica tra soggetti essenziali e importanti,
misure di esecuzione e sanzioni. Chiude con il caso del disservizio globale di \textbf{Microsoft~365}
del 12--13/11/2026 (``un evento, tre fascicoli'') e presenta il \textbf{laboratorio~5} di venerdì
20/11, una simulazione di notifiche NIS2 e GDPR. Messaggio di fondo: la domanda giusta non è mai ``è
NIS2?'', ma ``quale soggetto deve quale notifica, a chi e da quando''.

\section{Contenuti}

\subsection{Strategia nazionale e organismi di cooperazione}
In Italia l'\textbf{ACN} è l'autorità di regolazione per la NIS2; il \textbf{CSIRT Italia} ne è la
divisione tecnica (minacce, incidenti, risposta). La \textbf{Strategia Nazionale di Cybersicurezza
2022--2026} comprende \textbf{82 misure}, rivolte soprattutto alla pubblica amministrazione, con ricadute
su imprese e cittadini; sarà rinnovata a breve. Sull'attuazione della strategia l'ACN ha un ruolo più
politico che operativo: dà indirizzi ma non ha poteri per imporla.

Le slide elencano i compiti dei \textbf{CSIRT nazionali}: monitoraggio e analisi di minacce,
vulnerabilità e incidenti; preallarmi e allerte ai soggetti; risposta agli incidenti e assistenza;
analisi forense e consapevolezza situazionale; su richiesta, \textbf{scansione proattiva} dei sistemi di
un soggetto; partecipazione alla rete dei CSIRT. Sui livelli europei:
\begin{itemize}
  \item il \textbf{Cooperation Group} (Stati, Commissione, ENISA) è un organo consultivo, non di
    esecuzione: linee guida alle autorità, buone pratiche, \textbf{valutazioni coordinate dei rischi
    delle catene di approvvigionamento critiche}, indirizzi alla rete dei CSIRT e a EU-CyCLONe;
  \item \textbf{EU-CyCLONe} (autorità di gestione delle crisi degli Stati, più la Commissione negli
    incidenti su larga scala) si attiva per gli incidenti e le crisi su vasta scala: preparazione,
    valutazione dell'impatto, coordinamento e supporto alle decisioni a livello politico.
\end{itemize}
Immagine del docente: la rete dei CSIRT è il ``piano terra'' (lavoro quotidiano, scambio di minacce e
segnalazioni); quando si sale al livello politico e strategico europeo entra in gioco EU-CyCLONe.

\subsection{Governance: l'organo di gestione}
Gli organi di gestione dei soggetti essenziali e importanti \textbf{approvano} le misure di gestione del
rischio e ne \textbf{sorvegliano l'attuazione}; devono seguire una formazione e incoraggiare quella dei
dipendenti (slide). Il docente sottolinea la novità: di regola in Europa (diversamente dagli Stati
Uniti) la sanzione colpisce la società, non i consiglieri; con la NIS2 il componente dell'organo che non
fornisce indirizzo e sorveglianza sulla cybersicurezza risponde \textbf{personalmente}. Il consiglio
non deve gestire i singoli controlli, ma dare strategia, obiettivi e supervisione; questo non si
delega. È uno dei motivi per cui la NIS2 viene presa sul serio.

\subsection{Misure proporzionate e perimetro}
Le misure sono \textbf{tecniche, operative e organizzative}, basate su un approccio
\textbf{multi-rischio} (\emph{all-hazards}) e \textbf{proporzionate} all'esposizione al rischio, alla
dimensione e alla probabilità e gravità degli incidenti; devono prevenire o ridurre l'impatto degli
incidenti sui destinatari dei servizi e su altri servizi. Le slide riportano le dieci misure minime
(Sezione~\ref{sec:nis2-misure}). ``Proporzionato significa tutto e niente'': in ottica di
accountability bisogna dimostrare di aver valutato il rischio e l'efficacia dei controlli. Per il
docente il minimo indispensabile è: gestione degli incidenti con responsabilità chiare, politiche di
sicurezza, controllo e audit della supply chain (anche acquisti di hardware e infrastrutture), gestione
delle vulnerabilità e penetration test, valutazione dell'efficacia delle misure; MFA e comunicazioni
protette sono ``fortemente consigliate'', ove opportuno.

\paragraph{Il perimetro.} Un soggetto può rientrare nella NIS2 solo per alcuni prodotti o servizi.
Valutazioni, livelli di servizio e misure si fanno \textbf{solo sul perimetro NIS2}: ciò che è fuori
ambito resta fuori. Il difficile è che processi fuori perimetro possono incidere ``a cascata'' su
processi in perimetro: bisogna risalire fino al processo più granulare.

\paragraph{Il fornitore in crisi.} Nel disservizio di Microsoft, Microsoft ha i suoi obblighi di
notifica, ma anche ogni soggetto NIS2 che usa Outlook per erogare i propri servizi deve valutare se
notificare: l'incidente nella sua supply chain lo rende responsabile dei propri obblighi.

\subsection{Divulgazione coordinata delle vulnerabilità}
Dalle slide: ogni Stato designa un CSIRT \textbf{coordinatore}, intermediario di fiducia tra chi segnala
una vulnerabilità e il fabbricante o fornitore; chiunque può segnalare, anche in forma anonima;
per le vulnerabilità che toccano più Stati i coordinatori cooperano nella rete dei CSIRT; ENISA
gestisce una \textbf{banca dati europea delle vulnerabilità}, alimentata su base volontaria. Il docente
segnala il problema aperto dei \textbf{ricercatori di sicurezza}: alcune loro attività sono reato in
Italia e in molti Paesi europei; se i ricercatori migliori temono di essere incriminati e non
segnalano, l'ecosistema di scambio di informazioni costruito dalla NIS2 ne soffre. Se ne discute a
livello europeo.

\subsection{Incidente e incidente significativo}
L'\textbf{incidente} (art.~6, punto~6) è un evento che compromette disponibilità, \textbf{autenticità},
integrità o riservatezza di dati conservati, trasmessi o elaborati, o dei servizi offerti da reti e
sistemi informativi o accessibili attraverso di essi. Non ogni incidente si notifica: solo quello
\textbf{significativo}, cioè che ha causato o \textbf{può causare} una grave perturbazione operativa o
perdite finanziarie per il soggetto, oppure si è ripercosso o può ripercuotersi su altre persone
causando danni materiali o immateriali considerevoli.

\paragraph{La soglia nella pratica.} Secondo le linee guida ACN il soggetto definisce i propri
\textbf{livelli di servizio attesi} (una \emph{baseline}): superarli è un incidente significativo.
Esempio del docente: se il livello di servizio fissato è un'indisponibilità massima di 17 ore, un fermo
più lungo è significativo, uno più breve può non esserlo. Va però valutato anche l'effetto su terzi che
dipendono dai servizi del soggetto.

\paragraph{Materiale e immateriale.} Nell'inglese giuridico \emph{material} significa spesso
``rilevante''; nella NIS2 \emph{material / non-material} significa \textbf{patrimoniale / non
patrimoniale}: danno misurabile in denaro da un lato, danno alla reputazione o alla salute dall'altro
(si pensi alla sanità). Anche il danno non patrimoniale si può liquidare in giudizio.

\subsection{Chi notifica, a chi, quando}
Notificano i \textbf{soggetti} essenziali e importanti (le società, non i consiglieri) al
\textbf{CSIRT} competente; in Italia al CSIRT Italia tramite il portale ACN, non ``all'ACN'' in
generale. Le slide:
\begin{itemize}
  \item \textbf{24 ore} dalla consapevolezza: \emph{early warning}, indicando se si sospetta un atto
    illegittimo o malevolo e se l'impatto può essere transfrontaliero. Vale anche nel fine settimana.
    Contiene poche informazioni (per il docente anche meno della prima notifica GDPR): che cosa
    succede, quando, come è stato scoperto ed eventuali indicatori di compromissione;
  \item \textbf{72 ore}: notifica dell'incidente con valutazione iniziale di gravità e impatto e, ove
    disponibili, gli \textbf{indicatori di compromissione} (prestatori di servizi fiduciari: sempre
    24 ore);
  \item \textbf{su richiesta} del CSIRT: relazione intermedia. Il CSIRT raccoglie informazioni
    attivamente, non aspetta;
  \item \textbf{entro un mese dalla notifica dell'incidente}: relazione finale con descrizione
    dettagliata, tipo di minaccia o causa di fondo, misure di mitigazione, impatto transfrontaliero;
  \item se l'incidente è \textbf{ancora in corso}: relazione sui progressi a quella data e relazione
    finale entro un mese dalla gestione dell'incidente.
\end{itemize}
Il docente aggiunge che nella relazione finale vanno anche \emph{gap analysis} e piano di rimedio, e che
oltre i due mesi dall'incidente senza relazione finale ci si deve aspettare un audit di ACN e CSIRT.
Diversamente dal GDPR, dove l'integrazione della notifica non ha un termine finale, qui la mancata o
insufficiente relazione finale è di per sé una violazione.

\paragraph{Due orologi.} L'orologio NIS2 parte dalla consapevolezza dell'incidente significativo;
quello del GDPR (72 ore) dalla consapevolezza della violazione di dati personali, che può arrivare più
tardi, anche dopo la scadenza delle 72 ore NIS2. Le scadenze possono coincidere, ma non è automatico:
``avvio l'orologio NIS2, quindi tra 72 ore faccio entrambe le notifiche'' è un errore.

\paragraph{Dopo la notifica.} Il CSIRT risponde, ove possibile entro 24 ore dal preallarme, con un
primo riscontro e, su richiesta, indicazioni sulla mitigazione; può sapere cose che il soggetto ignora
(campagne in corso in altri Stati). Se l'incidente riguarda due o più Stati informa senza indebito
ritardo gli altri Stati ed ENISA, tutelando la sicurezza, gli interessi commerciali e la riservatezza
del soggetto. Il docente insiste sulla riservatezza: una mappa dell'infrastruttura di un gruppo
multinazionale, se divulgata, diventa essa stessa un rischio.

\begin{esame}
\begin{itemize}
  \item Orologi NIS2: \textbf{24 h / 72 h / 1 mese} dalla \emph{notifica} dell'incidente, al CSIRT
    (in Italia CSIRT Italia, tramite il portale ACN). Gli orologi NIS2 e GDPR sono \textbf{distinti}.
  \item Responsabilità \textbf{personale} dei componenti dell'organo di gestione.
  \item Vigilanza asimmetrica: essenziali \emph{ex ante}, importanti \emph{ex post}.
  \item Sanzioni 10~milioni / 2\% (essenziali) e 7~milioni / 1,4\% (importanti); una sola sanzione per
    la stessa condotta tra NIS2 e GDPR.
\end{itemize}
\end{esame}

\subsection{Vigilanza: essenziali e importanti}
Tabella delle slide, commentata dal docente:
\begin{center}
\small
\begin{tabularx}{\linewidth}{>{\raggedright\arraybackslash}p{2.8cm} X X}
\toprule
 & \textbf{Essenziali} & \textbf{Importanti}\\
\midrule
Approccio & \emph{ex ante}, basato sul rischio & \emph{ex post}, su indizi di non conformità\\
Ispezioni & in loco e a distanza, anche a campione & in loco e a distanza, \emph{ex post}\\
Audit & periodici e mirati, più \emph{ad hoc} (es.\ dopo un incidente significativo) & mirati\\
Scansioni & su criteri oggettivi, non discriminatori, equi e trasparenti & stessi criteri\\
Informazioni & richieste di informazioni, dati, documenti, prove di attuazione & le stesse, per la
  valutazione \emph{ex post}\\
Costo dell'audit indipendente & a carico del soggetto & a carico del soggetto\\
\bottomrule
\end{tabularx}
\end{center}
Commento del docente: un soggetto essenziale mette a budget gli audit; un importante che non subisce
incidenti ``non riceverà visite'', ma se notifica può essere sottoposto a verifica e pagarne il costo,
il che crea un incentivo (sbagliato) a non notificare.

\subsection{Esecuzione e sanzioni}
Scala delle misure (slide e docente): avvertimenti e \textbf{istruzioni vincolanti} (anche con una
\emph{gap analysis} e le misure da adottare); ordini di cessare la condotta, di conformarsi
all'art.~21 o agli obblighi di notifica entro un termine, di informare i destinatari esposti a una
minaccia significativa, di attuare le raccomandazioni di un audit; nomina di un responsabile del
monitoraggio e pubblicazione della violazione. Se le misure restano inefficaci, l'autorità fissa un
termine e poi può \textbf{sospendere} certificazioni o autorizzazioni relative ai servizi e chiedere il
\textbf{divieto temporaneo di funzioni dirigenziali} fino al livello di amministratore delegato o
rappresentante legale: misure di ultima istanza, con garanzie procedurali. Secondo il docente la leva
sulle persone (sospensione, stipendio) spesso funziona meglio di quella sulla società.

\paragraph{Sanzioni pecuniarie.} Per la violazione delle misure di sicurezza e degli obblighi di
notifica: massimo di almeno \textbf{10 milioni di euro o 2\%} del fatturato mondiale annuo, se
superiore, per gli \textbf{essenziali} (il livello più basso della scala GDPR); almeno \textbf{7
milioni o 1,4\%} per gli \textbf{importanti}. Gli Stati possono prevedere penalità di mora. Secondo il
docente le sanzioni sono \textbf{triplicate} (30 milioni o 6\%) per chi non si registra nel registro
nazionale dei soggetti NIS2 o non comunica il punto di contatto: il regolatore vuole prima di tutto
sapere chi vigilare, e molte imprese hanno provato a non registrarsi.

\subsection{NIS2 e gli altri strumenti}
Tabella delle slide: la NIS2 si applica fatti salvi GDPR, direttiva ePrivacy, le direttive 2011/93/UE e
2013/40/UE e la direttiva CER. Le \textbf{entità critiche} CER sono essenziali a prescindere dalla
dimensione; il \textbf{DORA} è \emph{lex specialis} (un'entità finanziaria in ambito di entrambe applica
solo il DORA); il \textbf{CRA} riguarda la sicurezza dei prodotti ed è complementare agli obblighi
organizzativi NIS2; la NIS1 (Dir.~2016/1148) è abrogata dal 18/10/2024.

\paragraph{NIS2 e GDPR.} Le misure di sicurezza si sovrappongono (art.~32 GDPR). Se una violazione NIS2
può comportare un data breach da notificare ai sensi dell'art.~33 GDPR, l'autorità NIS informa senza
indebito ritardo il Garante; se il Garante ha già inflitto una sanzione ai sensi dell'art.~58(2)(i)
GDPR, l'autorità NIS \textbf{non} può sanzionare la stessa condotta ai sensi dell'art.~34 NIS2, ma può
adottare le misure di esecuzione. ``Una sola sanzione'' (\emph{ne bis in idem}): per il docente è anche
il motivo per cui gli importi sono allineati. Le due notifiche corrono in parallelo con termini
diversi; gli Stati sono incoraggiati a creare un \textbf{punto di ingresso unico} per le notifiche.

\subsection{Messaggi chiave (slide 36)}
\begin{itemize}
  \item La NIS2 (Dir.~(UE) 2022/2555) è un quadro unico per \textbf{18 settori critici}; allarga molto
    l'ambito (\textbf{da 8 a 25 settori}, secondo la slide) e sostituisce le categorie della NIS1 con
    soggetti essenziali e importanti; l'elenco è tenuto a livello nazionale e riportato nel registro
    ENISA.
  \item Gli obblighi di sicurezza sono minimi, basati sul rischio e multi-rischio (almeno dieci misure),
    con divulgazione coordinata delle vulnerabilità e banca dati europea.
  \item Notifica degli incidenti significativi 24 h / 72 h / 1 mese; vigilanza asimmetrica; in Italia
    il D.lgs.~138/2024 designa l'ACN (con il CSIRT Italia al suo interno) e fa decorrere i termini dalla
    comunicazione di inserimento nell'elenco NIS.
\end{itemize}
Il docente aggiunge che tra i settori ci sono anche \textbf{manifatture}, perché i sistemi informativi
sono incorporati nella produzione, e che l'ACN ha per la NIS2 il ruolo che il Garante ha per la
protezione dei dati.

\subsection{Caso di studio: il disservizio globale di Microsoft~365}
Il docente lascia le slide da leggere e ne riassume il punto. Dalle slide:
\begin{itemize}
  \item \textbf{Fatti.} Il 12--13/11/2026 Outlook, Teams, OneDrive, SharePoint, le app Office e
    Copilot sono stati degradati o non disponibili per buona parte di una giornata, con problemi in
    almeno 74 Paesi e ripercussioni sui servizi di terzi integrati con Microsoft~365. Causa (rapporto
    preliminare del 15/11): un difetto software introdotto da un aggiornamento di luglio 2026 al
    gateway di rete, scatenato da un picco di carico mentre circa il 14\% della capacità era fuori
    servizio, propagato a cascata (al minimo oltre l'80\% dei server non rispondeva). Nessuna prova di
    attività malevola: per la definizione di incidente \textbf{non serve un atto malevolo}.
  \item \textbf{Il fascicolo del fornitore.} Un fornitore di cloud è nell'Allegato~I (infrastrutture
    digitali): sopra le soglie della media impresa è essenziale. Per i fornitori di cloud la
    significatività non dipende da un livello di servizio auto-dichiarato: l'\textbf{art.~7 del
    Regolamento di esecuzione (UE) 2024/2690} fissa soglie (servizio del tutto indisponibile per più
    di 30 minuti; disponibilità limitata per più di un'ora per oltre il 5\% degli utenti nell'Unione o
    oltre un milione di utenti, se inferiore; compromissione di integrità, riservatezza o autenticità
    per sospetta azione malevola sopra la stessa soglia). Solo il terzo criterio richiede la malizia:
    ``nessuna traccia di attacco'' non esclude nulla. La notifica va al CSIRT dello Stato dello
    \textbf{stabilimento principale} (art.~26), dove si prendono le decisioni di gestione del rischio.
    Tre doveri distinti, non assolti da una pagina di stato pubblica: notificare il CSIRT, informare i
    destinatari dei servizi colpiti (art.~23(1)), comunicare misure e rimedi per le minacce
    significative (art.~23(2)).
  \item \textbf{Un evento, tre fascicoli} per il cliente soggetto NIS2, che non è sollevato dalla
    notifica del fornitore: (1) l'incidente significativo del fornitore; (2) il \textbf{proprio}
    incidente significativo, valutato sui propri livelli di servizio attesi (tipologia ACN
    \textbf{IS-3}, superamento dei livelli definiti nella misura DE.CM-01) e sull'art.~23(3), da
    notificare al CSIRT Italia tramite il portale ACN con gli orologi 24 h / 72 h / 1 mese; (3) il
    fascicolo di rischio della \textbf{supply chain} (art.~21(2)(d) e~(3)): la dipendenza da un solo
    fornitore è essa stessa un rischio da documentare e monitorare.
  \item \textbf{Livello di servizio non è SLA}: uno SLA contrattuale violato, o assente, non decide
    l'IS-3; lo decidono i livelli di servizio attesi documentati dal soggetto, senza i quali l'IS-3 non
    si può nemmeno rilevare. Mentre l'evento è aperto vanno raccolti aggiornamenti e rapporto del
    fornitore, servizi e ore di indisponibilità, chi è stato colpito, soluzioni alternative e
    l'eventuale compromissione di dati (il binario GDPR corre in parallelo).
\end{itemize}

\begin{nota}[Il problema del ``livello'' irlandese]
Le slide segnalano che al 4/11/2026 l'Irlanda, dove Microsoft ha lo stabilimento principale europeo,
non aveva ancora recepito la NIS2 (parere motivato del 7/05/2025, deferimento alla Corte di giustizia
deciso l'8/07/2026; resta in vigore la normativa di attuazione della NIS1): se lo stabilimento principale
è in uno Stato che non ha recepito, l'obbligo del fornitore può non sorgere lì.
\end{nota}

\subsection{Il laboratorio~5 (venerdì 20/11, online)}
Il docente presenta il laboratorio: \textbf{EnergiaTech S.p.A.}, scenario fittizio di ransomware su
un distributore regionale di energia elettrica. Esercizio a squadre di \textbf{90 minuti}, autogestito
(niente facilitatore e niente ``imprevisti'' durante l'esercizio).
\begin{itemize}
  \item \textbf{Profilo} (slide): soggetto \textbf{essenziale} (``non dovete chiedervi se è
    importante''), 1.200 dipendenti, 15 sottostazioni e 3 centri di distribuzione, 450.000 utenti (non
    tutti italiani: ``un indizio'' per il GDPR). \textbf{Attacco}: spear phishing mirato, LockBit~3.0
    con doppia estorsione, 850 endpoint cifrati, 487~GB esfiltrati, richiesta di 50~BTC con minaccia di
    pubblicazione in 7 giorni, permanenza non rilevata di circa 30 giorni. La distribuzione elettrica
    non si è mai interrotta (OT separata dall'IT); fatturazione, CRM, posta e assistenza clienti
    bloccati per giorni.
  \item \textbf{Che cosa si valuta}: il processo, non l'analisi forense (i fatti tecnici sono dati
    nel dossier, consegnato solo venerdì): classificare (incidente significativo ex art.~23(3)?
    quale tipologia ACN, da IS-1 a IS-4?), \textbf{fissare l'inizio degli orologi} e motivarlo, scrivere
    in modo difendibile, tenere separati i binari NIS2 (CSIRT/ACN) e GDPR (Garante), \textbf{non
    inventare mai un fatto} (ciò che non si sa resta un segnaposto dichiarato, oppure si chiede).
  \item \textbf{Consegne}: \emph{early warning} (10 campi), notifica dell'incidente (campi 11--19 più
    allegato IoC), relazione finale (le quattro voci dell'art.~23(4)(d), più piano correttivo e
    \emph{lessons learned}; riferimento nazionale art.~25 D.lgs.~138/2024), notifica ex art.~33 GDPR al
    Garante (modulo A--O), nota di accompagnamento con le risposte alle domande di autovalutazione; un
    file ZIP per squadra su SharePoint. File in markdown o PDF, nessuna formalità. Nessuna relazione
    intermedia; la comunicazione agli interessati (art.~34) non è richiesta, ma una motivazione solida
    del perché farla o non farla può far guadagnare punti.
  \item \textbf{Canale}: portale ACN / CSIRT Italia, non ``l'ACN'' in generale: il docente avverte che
    sbagliarlo costa molto nel voto.
  \item \textbf{Materiale}: dopo la lezione il docente carica le slide, un manuale preparatorio
    (\emph{pre-flight handbook}) e modelli compilati su un altro caso (SalusNova, piattaforma
    sanitaria), da usare come struttura, non come risposte. Ci si può discostare dal modello
    motivandolo.
\end{itemize}
Le slide elencano anche le trappole da controllare: decidere la significatività solo sui numeri,
sbagliare l'orologio, fondere NIS2 e GDPR in un documento, affermare ciò che il dossier non supporta,
dimenticare una notifica o comunicazione dovuta, lasciare segnaposto vuoti, trattare la notifica come
un'ammissione di responsabilità (non lo è, art.~23(1)).

\section{Avvisi}
\begin{itemize}
  \item Venerdì 20/11: laboratorio~5 (NIS2 + GDPR), da remoto. Il docente lo ha reso più semplice del
    laboratorio~4 (audit), giudicato troppo complesso.
  \item Più avanti potrebbe esserci un laboratorio più ampio su un incidente che coinvolge più norme
    (NIS2, GDPR, CRA, DORA, AI Act), forse con più di 90 minuti; per ora i docenti hanno deciso di non
    aggiungere il CRA a questo laboratorio per non creare confusione.
  \item Slide NIS2 aggiornate, slide del laboratorio, manuale preparatorio e modelli: caricati sulla
    piattaforma dopo la lezione.
\end{itemize}

\section{Domande aperte e punti da verificare}
\begin{verifica}
\begin{itemize}
  \item \textbf{Sanzioni triplicate} per mancata registrazione (30 milioni o 6\%): detto a voce, non
    nelle slide; da verificare nel D.lgs.~138/2024.
  \item \textbf{Relazione finale}: il docente dice sia ``un mese meno tre giorni dopo la seconda
    notifica'' sia ``un mese dall'inizio dell'incidente''; la slide e l'art.~23(4)(d) dicono un mese
    dalla \emph{notifica dell'incidente} (la 72 ore). Il limite dei ``due mesi'' con audit è
    un'indicazione del docente, non della norma.
  \item \textbf{IoC}: nelle slide sono nella notifica a 72 ore (``ove disponibili''); a voce il docente
    li cita già nel preallarme, se l'analisi li ha prodotti.
  \item \textbf{Settori}: ``18 settori critici'' e ``da 8 a 25 settori'' nella stessa slide; 18 = 11
    (Allegato~I) + 7 (Allegato~II); il 25 conterebbe anche i sottosettori o gli allegati italiani.
  \item \textbf{Strategia nazionale}: il docente parla di strategia ``di quattro anni (2022--2026)'' e
    ``triennale''; il titolo nelle slide è Strategia Nazionale di Cybersicurezza 2022--2026.
  \item Il docente cita una sanzione del Garante a Meta ``lo scorso fine settimana'' (direttiva
    ePrivacy): non approfondita.
  \item \textbf{Docenti}: Gilardi non si presenta; il secondo intervento (sulla proposta di aggiungere il
    CRA al laboratorio) è attribuito a Vaciago dal contesto (``ne abbiamo parlato con Giuseppe'').
  \item Lo scenario del laboratorio è datato ``domenica 16 novembre 2025'' (Day~0) nelle slide.
\end{itemize}
\end{verifica}
